\section{Additional Method Details}
The main paper describes how domains are connected and how source actions are retained. This section records an optional reduction for a fixed world and the generic search implementation used in earlier development studies. The learned-domain evaluation uses the frozen groups and Glaive settings reported in the main paper.

\subsection{Optional Reduction for a Fixed World}
For a supplied world, let $U_i$ contain the objects satisfying source $i$'s immutable root-membership predicate. Every declared source constant must retain its recorded memberships, and every positive subtype membership must imply root membership. This admission check matters because a constant literal in an effect has no action-parameter guard. The selector rejects missing constant memberships rather than inferring or restoring facts.

Construct a graph over sources. Sources sharing a target predicate are adjacent if that predicate is nullary or their root universes overlap. A positive-arity interface cannot communicate a tuple between disjoint root universes, regardless of its argument permutation. Let $S$ be the union of graph components containing every declared owner of a positive or negative goal atom. The reduction retains all actions from $S$, together with every original declaration, constant, and the exact mapped problem. It does not recompile a smaller object universe. Its certificate records the source universes and retained owners; reuse with changed universes or an unsupported goal owner requires recertification.

\begin{proposition}[Fixed-universe reduction]
Assume the stated membership compilation, immutable memberships covering source constants, source-owned action variants without additional adapters, and a ground goal whose atoms lie in their declared source universes. Let $D_S$ remove only the actions outside $S$. Then the supplied goal is reachable in $D_S$ if and only if it is reachable in the full compiled domain.
\end{proposition}
\begin{proof}
Membership-guarded parameters and quantifiers, together with constant coverage, restrict each source's effective predicate reads and writes to tuples from its own $U_i$. A nullary shared predicate creates an edge. For a positive-arity shared predicate, an excluded source and a retained source have disjoint universes, so their possible tuples are disjoint. Excluded actions therefore leave all retained source-tuple truth values unchanged. The truth of retained conditions and their conditional effects depends only on those tuples and immutable memberships; declarations and quantifier universes remain identical. Deleting excluded actions from a full solution consequently preserves each retained action's applicability and the supplied goal. Conversely, every reduced solution is a full-domain solution because retained actions and the problem are unchanged.
\end{proof}

\paragraph{Boundary.}
The statement concerns goal reachability, not full-state trace equivalence: excluded independent actions can change other tuples or events. It provides no minimum-component, bounded-search success, narrative-quality, or interface-semantic guarantee, and it does not apply to initial-only outcome ranking. The proof argument and its relationship to existing modularity and decomposition work still require independent technical review. A deliberately malformed mapped-world control exposes a lost producer when constant coverage is omitted; the strengthened admission check rejects that input. Rechecking all thirty valid component inputs in the matched source-world block leaves their generated domains byte-identical.

\subsection{Generic Search and Outcome Generation}
\paragraph{Goal-directed mode.}
The default reference search uses bounded BFS and static-positive parameter projection. The latter restricts candidate domains only with conjunctive static Boolean atoms having false defaults; remaining correlations and dynamic conditions are checked by the simulator. State, depth, grounding, and subprocess time limits are reported distinctly. The optional Pyperplan reference uses greedy best-first search with hFF after recorded UP compilations; its adapter enforces only a subprocess time budget. Every returned plan is independently checked in the original exported model after compiler plan lifting.

\paragraph{Initial-only mode.}
Without a goal, bounded BFS returns the first distinct noninitial states in discovery order. Each candidate records its positive and negative changes from $I$, a nonempty plan, and validation against the corresponding generated endpoint condition. This is a declared outcome-generation policy, not narrative ranking or success against independently specified goals. Resource cutoffs remain visible even when some outcomes have been generated.

For false-default Boolean states, the implementation maintains the complete set of true ground atoms. Conjunctive positive preconditions project parameter choices against this current-state index; remaining conditions are checked by simulation. All expanded effect targets update the index from the successor state, including deletions and conditional or quantified effects. Endpoint differences are the symmetric difference of the initial and successor true-atom sets. Grounding and fluent budgets count unique visited projected bindings and effect targets on visited state-changing transitions, respectively. These counters differ from eager enumeration. Authored finite controls agree with unpruned Cartesian BFS on result order, plans, and complete changes; they do not establish corpus-wide scaling.
